%----------------------------------------------------------------------------------------
% EXPONENTIËLE FUNCTIES. Alle opgaven zijn eigen werk.
%
% Opbouw: Start, leerstof (theorie, werkwijze of uitleg, voorbeeld, korte oefening), en op een
% nieuwe pagina de Opdrachten (door elkaar, stijgende moeilijkheid), Koppel en Test jezelf.
% Labels van dit hoofdstuk beginnen met h1-.
%----------------------------------------------------------------------------------------

\mychapter{1}{Exponentiële functies}

\lkr{Leerplan B+S': LPD 7-12 (exponentiële en logaritmische functies), het deel over machten met rationale exponenten, de exponentiële functie en exponentiële groei. Vergelijkingen waarbij een logaritme nodig is, komen in hoofdstuk 2: hier lossen leerlingen alleen vergelijkingen op waarin beide leden een macht van hetzelfde grondtal zijn. De transformaties van grafieken komen uit de cursus Algemene functies.}

Een blad papier is ongeveer $0{,}1$ mm dik. Vouw je het dubbel, dan verdubbelt de dikte, en bij elke volgende vouw opnieuw. Na tien vouwen is het pakje al meer dan tien centimeter dik. Zo'n groei, waarbij elke stap met dezelfde \textit{factor} vermenigvuldigt, heet exponentieel. Je vindt ze bij intrest, bevolkingen, bacteriën en radioactief verval.

\begin{oefeningen}[Start]
\oef{h1-start-1}

\oef{h1-verken}
\end{oefeningen}

%========================================================================================
\section{Machten met rationale exponenten}

Uit de 2de graad ken je machten met gehele exponenten, zoals $2^3 = 8$, $2^0 = 1$ en $2^{-3} = \frac18$, en de rekenregels om ermee te rekenen. Voor exponentiële groei hebben we meer nodig. Verdubbelt iets elk jaar, hoeveel keer zo groot is het dan na een half jaar? Dat is $2^{\frac12}$, en zo'n exponent heeft voorlopig geen betekenis. In deze sectie geven we ze een betekenis, en tonen we aan dat de vertrouwde rekenregels blijven gelden.

\subsection{Gebroken exponenten}

Hoe moet $8^{\frac13}$ heten? Als de rekenregel $\bigl(a^r\bigr)^s = a^{rs}$ ook voor breuken moet gelden, dan is
\[\bigl(8^{\frac13}\bigr)^3 = 8^{\frac13 \cdot 3} = 8^1 = 8 .\]
Dus is $8^{\frac13}$ een getal waarvan de derde macht $8$ is: de derdemachtswortel $\sqrt[3]{8} = 2$. Zo komen we tot de definitie.

\begin{theorie}
Voor $a > 0$, $m \in \mathbb{Z}$ en $n \in \mathbb{N} \setminus \{0\}$ is
\[\formule{a^{\frac{m}{n}} = \sqrt[n]{a^m}}.\]
\end{theorie}

De noemer van de exponent is dus de wortel, de teller de macht. Een $n$-de machtswortel $\sqrt[n]{b}$ van een positief getal $b$ is het positieve getal waarvan de $n$-de macht $b$ is. Er is maar één zo'n positief getal. Die eigenschap gebruiken we in elk bewijs van deze sectie:
\[c > 0 \ \text{ en } \ c^n = a^m \quad \Longrightarrow \quad c = \sqrt[n]{a^m} = a^{\frac{m}{n}} .\]
Daarmee zie je ook dat het niet uitmaakt hoe je de breuk schrijft: $\frac13 = \frac26$, en inderdaad is $\sqrt[6]{8^2} = \sqrt[3]{8}$, want $\bigl(\sqrt[3]{8}\bigr)^6 = \bigl(\bigl(\sqrt[3]{8}\bigr)^3\bigr)^2 = 8^2$.

Waarom $a > 0$? Met een negatief grondtal loopt het mis: $(-8)^{\frac13}$ zou $\sqrt[3]{-8} = -2$ zijn, maar $(-8)^{\frac26} = \sqrt[6]{(-8)^2} = \sqrt[6]{64} = 2$, terwijl $\frac13 = \frac26$. En $(-4)^{\frac12} = \sqrt{-4}$ bestaat niet. Daarom spreken we af dat het grondtal positief is.

\voorbeeld{}{Bereken zonder rekentoestel: $8^{\frac23}$.\label{vb:h1-macht}}{%
\uitleg{De noemer van de exponent is de wortel, de teller de macht. Neem eerst de wortel: dan blijven de getallen klein.}
\begin{stappenbreed}
8^{\frac23} = \sqrt[3]{8^2} = \bigl(\sqrt[3]{8}\bigr)^2 & De wortel van een macht is de macht van de wortel. \\
= 2^2 = 4 & Controle: $4^3 = 64 = 8^2$.
\end{stappenbreed}}

\begin{oefeningen}
\oef{h1-k-macht}
\end{oefeningen}

\subsection{Product van machten met hetzelfde grondtal}

Voor gehele exponenten weet je dat $a^3 \cdot a^4 = a^7$: je telt de exponenten op. Geldt dat ook voor breuken, nu $a^{\frac{m}{n}}$ een wortel is? Dat is niet vanzelfsprekend, en daarom bewijzen we het.

\begin{theorie}
\begin{stelling}{Product van machten}
Voor $a > 0$ en rationale $r$ en $s$ is \qquad $\formule{a^r \cdot a^s = a^{r + s}}$
\end{stelling}
\begin{bewijs}
Schrijf $r$ en $s$ op dezelfde noemer: $r = \frac{p}{n}$ en $s = \frac{q}{n}$, met $p, q \in \mathbb{Z}$ en $n \geqslant 1$. Stel $c = a^r \cdot a^s$, een positief getal. Met de rekenregels voor gehele exponenten is
\bewijsregel{c^n = \bigl(a^r\bigr)^n \cdot \bigl(a^s\bigr)^n = a^p \cdot a^q = a^{p + q},}
want $a^r = \sqrt[n]{a^p}$ heeft als $n$-de macht $a^p$, en net zo voor $a^s$. Dus is $c$ de $n$-de machtswortel van $a^{p+q}$:
\bewijsregel{c = \sqrt[n]{a^{p + q}} = a^{\frac{p + q}{n}} = a^{r + s}.}
\end{bewijs}
\end{theorie}

Bij het vermenigvuldigen van machten met hetzelfde grondtal tel je dus de exponenten op, ook als het breuken zijn.

\voorbeeld{}{Bereken zonder rekentoestel: $3^{\frac34} \cdot 3^{\frac54}$.\label{vb:h1-product}}{%
\uitleg{Elke factor apart is een vierdemachtswortel die je niet uit het hoofd kent. Samen wordt het een gehele exponent.}
\begin{stappenbreed}
3^{\frac34} \cdot 3^{\frac54} = 3^{\frac34 + \frac54} & Zelfde grondtal: exponenten optellen. \\
= 3^{2} = 9 & $\frac34 + \frac54 = \frac84 = 2$.
\end{stappenbreed}}

\begin{oefeningen}
\oef{h1-k-product}
\end{oefeningen}

\subsection{Quotiënt van machten met hetzelfde grondtal}

Delen is het omgekeerde van vermenigvuldigen. Daarom volgt de regel voor een quotiënt meteen uit die voor een product.

\begin{theorie}
\begin{stelling}{Quotiënt van machten}
Voor $a > 0$ en rationale $r$ en $s$ is \qquad $\formule{\frac{a^r}{a^s} = a^{r - s}}$ \qquad en in het bijzonder \qquad $\formule{a^{-s} = \frac{1}{a^s}}$
\end{stelling}
\begin{bewijs}
Met de productregel is
\bewijsregel{a^{r - s} \cdot a^s = a^{(r - s) + s} = a^r .}
Deel beide leden door $a^s$, dat positief is: $a^{r - s} = \dfrac{a^r}{a^s}$. Neem je $r = 0$, dan staat er
\bewijsregel{a^{-s} = \frac{a^0}{a^s} = \frac{1}{a^s} .}
\end{bewijs}
\end{theorie}

Bij delen trek je de exponenten af. Een negatieve exponent betekent: één gedeeld door de macht met de positieve exponent.

\voorbeeld{}{Schrijf als één macht van $x$ (met $x > 0$): $\dfrac{\sqrt[3]{x^2}}{\sqrt{x}}$.\label{vb:h1-quotient}}{%
\uitleg{Schrijf eerst elke wortel als macht. Dan staat er een quotiënt van machten met hetzelfde grondtal.}
\begin{stappenbreed}
\frac{\sqrt[3]{x^2}}{\sqrt{x}} = \frac{x^{\frac23}}{x^{\frac12}} & Wortels als machten. \\
= x^{\frac23 - \frac12} = x^{\frac16} & Exponenten aftrekken: $\frac46 - \frac36 = \frac16$.
\end{stappenbreed}}

\begin{oefeningen}
\oef{h1-k-quotient}
\end{oefeningen}

\subsection{Macht van een macht}

Met deze regel begon de sectie: we wilden dat $\bigl(8^{\frac13}\bigr)^3 = 8^1$. Nu tonen we aan dat hij voor alle rationale exponenten klopt.

\begin{theorie}
\begin{stelling}{Macht van een macht}
Voor $a > 0$ en rationale $r$ en $s$ is \qquad $\formule{\bigl(a^r\bigr)^s = a^{rs}}$
\end{stelling}
\begin{bewijs}
Schrijf $r = \frac{p}{n}$ en $s = \frac{q}{m}$, met $p, q \in \mathbb{Z}$ en $n, m \geqslant 1$. Stel $b = a^r$ en $c = b^s = \sqrt[m]{b^q}$, twee positieve getallen. Dan is $c^m = b^q$ en $b^n = a^p$, en met de rekenregels voor gehele exponenten
\bewijsregel{c^{mn} = \bigl(c^m\bigr)^n = \bigl(b^q\bigr)^n = \bigl(b^n\bigr)^q = \bigl(a^p\bigr)^q = a^{pq}.}
Dus is $c$ de $mn$-de machtswortel van $a^{pq}$:
\bewijsregel{c = \sqrt[mn]{a^{pq}} = a^{\frac{pq}{mn}} = a^{rs}.}
\end{bewijs}
\end{theorie}

Bij een macht van een macht vermenigvuldig je de exponenten. Zo is ook $\sqrt[3]{a^{\frac34}} = \bigl(a^{\frac34}\bigr)^{\frac13} = a^{\frac14}$: een wortel uit een macht deelt de exponent.

\voorbeeld{}{Bereken zonder rekentoestel: $\bigl(16^{\frac38}\bigr)^{\frac23}$.\label{vb:h1-machtmacht}}{%
\uitleg{$16^{\frac38}$ is een achtstemachtswortel. Reken die niet uit: vermenigvuldig eerst de exponenten.}
\begin{stappenbreed}
\bigl(16^{\frac38}\bigr)^{\frac23} = 16^{\frac38 \cdot \frac23} = 16^{\frac14} & Exponenten vermenigvuldigen. \\
= \sqrt[4]{16} = 2 & $2^4 = 16$.
\end{stappenbreed}}

\begin{oefeningen}
\oef{h1-k-machtmacht}
\end{oefeningen}

\needspace{14\baselineskip}
\subsection{Macht van een product en van een quotiënt}

De vorige drie regels gingen over één grondtal met verschillende exponenten. De laatste twee gaan over verschillende grondtallen met dezelfde exponent.

\begin{theorie}
\begin{stelling}{Macht van een product}
Voor $a, b > 0$ en rationale $r$ is \qquad $\formule{(ab)^r = a^r \cdot b^r}$
\end{stelling}
\begin{bewijs}
Schrijf $r = \frac{p}{n}$, met $p \in \mathbb{Z}$ en $n \geqslant 1$. Stel $c = a^r \cdot b^r$, een positief getal. Dan is
\bewijsregel{c^n = \bigl(a^r\bigr)^n \cdot \bigl(b^r\bigr)^n = a^p \cdot b^p = (ab)^p ,}
dus $c = \sqrt[n]{(ab)^p} = (ab)^r$.
\end{bewijs}
\end{theorie}

Je mag de exponent dus over de factoren verdelen, of omgekeerd: twee machten met dezelfde exponent samennemen.

\voorbeeld{}{Bereken zonder rekentoestel: $2^{\frac13} \cdot 4^{\frac13}$.\label{vb:h1-prodgrondtal}}{%
\uitleg{De grondtallen zijn verschillend, dus de exponenten optellen kan niet. Maar de exponenten zijn gelijk: neem de grondtallen samen.}
\begin{stappenbreed}
2^{\frac13} \cdot 4^{\frac13} = (2 \cdot 4)^{\frac13} & Zelfde exponent: grondtallen vermenigvuldigen. \\
= 8^{\frac13} = \sqrt[3]{8} = 2 &
\end{stappenbreed}}

\begin{oefeningen}
\oef{h1-k-prodgrondtal}
\end{oefeningen}

De regel voor een quotiënt volgt weer uit die voor een product.

\begin{theorie}
\begin{stelling}{Macht van een quotiënt}
Voor $a, b > 0$ en rationale $r$ is \qquad $\formule{\Bigl(\frac{a}{b}\Bigr)^r = \frac{a^r}{b^r}}$
\end{stelling}
\begin{bewijs}
Met de vorige regel is
\bewijsregel{\Bigl(\frac{a}{b}\Bigr)^r \cdot b^r = \Bigl(\frac{a}{b} \cdot b\Bigr)^r = a^r .}
Deel beide leden door $b^r$, dat positief is.
\end{bewijs}
\end{theorie}

\voorbeeld{}{Bereken zonder rekentoestel: $\dfrac{54^{\frac13}}{2^{\frac13}}$.\label{vb:h1-quotgrondtal}}{%
\uitleg{Geen van beide is een mooie derdemachtswortel, maar het quotiënt van de grondtallen wel.}
\begin{stappenbreed}
\frac{54^{\frac13}}{2^{\frac13}} = \Bigl(\frac{54}{2}\Bigr)^{\frac13} & Zelfde exponent: grondtallen delen. \\
= 27^{\frac13} = 3 &
\end{stappenbreed}}

\begin{oefeningen}
\oef{h1-k-quotgrondtal}
\end{oefeningen}

\subsection{De rekenregels samen}

In een echte opgave heb je vaak meerdere regels na elkaar nodig. Schrijf dan eerst alle wortels als machten.

\voorbeeld{}{Schrijf als één macht van $x$ (met $x > 0$): $\dfrac{\sqrt{x} \cdot \sqrt[3]{x}}{x^{-\frac16}}$.\label{vb:h1-rekenregels}}{%
\uitleg{Schrijf elke wortel als macht en tel de exponenten op of trek ze af.}
\begin{stappenbreed}
\frac{x^{\frac12} \cdot x^{\frac13}}{x^{-\frac16}} = x^{\frac12 + \frac13 - \left(-\frac16\right)} & Wortels als machten, dan product- en quotiëntregel. \\
= x^{1} = x & $\frac36 + \frac26 + \frac16 = 1$.
\end{stappenbreed}}

Pas op: er is geen regel voor een som. Zo is $2^{\frac12} + 2^{\frac12} = 2 \cdot 2^{\frac12} = 2^{\frac32} \approx 2{,}83$, en niet $2^{\frac12 + \frac12} = 2$.

\begin{oefeningen}
\oef{h1-machten}
\end{oefeningen}

%========================================================================================
\section{De exponentiële functie}

We kennen nu $a^x$ voor elke rationale $x$. Om er een functie op heel $\R$ van te maken, ontbreken nog de irrationale exponenten.

\subsection{Machten met een reële exponent}

Wat betekent $2^{\sqrt2}$? We kunnen $\sqrt2 = 1{,}414\,21\ldots$ steeds beter benaderen met breuken, en van die breuken kennen we de macht:

\figuren{\renewcommand{\arraystretch}{1.2}%
\begin{tabular}{c|ccccc}
$x$ & $1{,}4$ & $1{,}41$ & $1{,}414$ & $1{,}414\,2$ & $\sqrt2$ \\ \hline
$2^x$ & $2{,}639\,0$ & $2{,}657\,4$ & $2{,}664\,7$ & $2{,}665\,1$ & $2{,}665\,1\ldots$
\end{tabular}}

De machten naderen een vast getal, en dat getal noemen we $2^{\sqrt2}$. Zo krijgt $a^x$ een betekenis voor elke reële $x$. De rekenregels blijven dan gelden. Dat bewijzen we niet: daarvoor heb je limieten nodig, en dat valt buiten het leerplan.

\begin{theorie}
Voor een grondtal $a > 0$ met $a \neq 1$ is $\formule{f(x) = a^x}$ een \textbf{\textcolor{accent}{exponentiële functie}}.
\end{theorie}

Het grondtal $1$ sluiten we uit, omdat $1^x = 1$ een constante functie is. Een negatief grondtal kan niet, zoals je bij de gebroken exponenten zag.

\subsection{Stijgen en dalen}

Hoe groter de exponent, hoe groter de macht? Dat klopt bij $2^x$, maar bij $\bigl(\frac12\bigr)^x$ is het net omgekeerd: $\bigl(\frac12\bigr)^1 = \frac12$ en $\bigl(\frac12\bigr)^2 = \frac14$. Het grondtal beslist.

\begin{theorie}
\begin{stelling}{Stijgen en dalen}
Als $a > 1$, dan is $a^x$ \textbf{stijgend}: \qquad $\formule{p < q \Rightarrow a^p < a^q}$

Als $0 < a < 1$, dan is $a^x$ \textbf{dalend}: \qquad $\formule{p < q \Rightarrow a^p > a^q}$
\end{stelling}
\begin{bewijs}
We bewijzen het voor rationale $p$ en $q$. Neem eerst $a > 1$, en een rationale $t = \frac{m}{n} > 0$ met $m, n \geqslant 1$. Uit $a > 1$ volgt $a^m > 1$. Was $a^t \leqslant 1$, dan ook $\bigl(a^t\bigr)^n = a^m \leqslant 1$. Dat kan niet, dus
\bewijsregel{a^t > 1 \quad \text{voor elke rationale } t > 0 .}
Neem nu $t = q - p > 0$. Dan is
\bewijsregel{a^q = a^p \cdot a^{q - p} > a^p \cdot 1 = a^p ,}
want $a^p > 0$. Is $0 < a < 1$, dan is $b = \frac1a > 1$, en met de rekenregels is $a^x = b^{-x}$. Uit $p < q$ volgt $-q < -p$, en omdat $b^x$ stijgt, is $a^q = b^{-q} < b^{-p} = a^p$.
\end{bewijs}
\end{theorie}

Voor irrationale exponenten geldt de stelling ook, maar dat bewijzen we niet. Ze zegt: bij een grondtal groter dan 1 hoort bij een grotere exponent een grotere macht, bij een grondtal tussen 0 en 1 een kleinere.

\voorbeeld{}{Welk getal is het grootst: $0{,}8^{-1{,}5}$ of $0{,}8^{-1{,}2}$? Antwoord zonder rekentoestel.\label{vb:h1-stijgen}}{%
\uitleg{Beide machten hebben hetzelfde grondtal. Kijk of dat grondtal groter of kleiner is dan 1, en vergelijk dan de exponenten.}
\begin{stappenbreed}
0 < 0{,}8 < 1 & Het grondtal ligt tussen 0 en 1: $0{,}8^x$ is dalend. \\
-1{,}5 < -1{,}2 & Vergelijk de exponenten. \\
0{,}8^{-1{,}5} > 0{,}8^{-1{,}2} & Dalend: de kleinste exponent geeft de grootste macht.
\end{stappenbreed}
Controle met het rekentoestel: $0{,}8^{-1{,}5} \approx 1{,}398$ en $0{,}8^{-1{,}2} \approx 1{,}307$.}

\begin{oefeningen}
\oef{h1-k-vergelijken}
\end{oefeningen}

\subsection{De grafiek}

Een macht van een positief grondtal is altijd positief, dus de grafiek ligt helemaal boven de $x$-as. Ze gaat door $(0, 1)$, want $a^0 = 1$, en door $(1, a)$.

\begin{theorie}
\begin{stelling}{Grafiek van een exponentiële functie}
De functie $f(x) = a^x$ heeft
\[\formule{\dom f = \R} \qquad \formule{\ber f = \left]0, +\infty\right[}.\]
Haar grafiek gaat door $(0, 1)$ en $(1, a)$, en heeft de $x$-as als \textbf{\textcolor{accent}{horizontale asymptoot}}.
\end{stelling}
\end{theorie}

Dat er geen negatieve beelden zijn, zag je hierboven. Dat elk positief getal wel een beeld is, en dat de grafiek de $x$-as nadert, bewijzen we niet: daarvoor heb je continuïteit en limieten nodig.

\naast{%
Een asymptoot is een rechte waar de grafiek steeds dichter bij komt, zonder ze te raken: $2^x$ wordt nooit nul, maar voor $x = -10$ is het al kleiner dan $0{,}001$.

Omdat $\bigl(\frac12\bigr)^x = 2^{-x}$, heeft $\bigl(\frac12\bigr)^x$ in $x$ dezelfde waarde als $2^x$ in $-x$. De grafiek van $\bigl(\frac12\bigr)^x$ is dus het spiegelbeeld van die van $2^x$ om de $y$-as. Hoe groter het grondtal, hoe steiler de grafiek rechts van de $y$-as.}{%
\begin{tikzpicture}[x=.55cm,y=.55cm,font=\scriptsize]
  \rooster{-3}{1}{3}{-1}{1}{5}
  \kromme{-3}{2.32}{2^(\x)}
  \kromme[dashed]{-2.32}{3}{2^(-\x)}
  \kromme[densely dotted]{-3}{1.46}{3^(\x)}
  \vol{(0,1)}
  \node[right,text=accent] at (2.2,4.6) {$2^x$};
  \node[left,text=accent] at (-2.1,4.6) {$\bigl(\frac12\bigr)^x$};
  \node[left,text=accent] at (1.3,4.9) {$3^x$};
\end{tikzpicture}}

Met de grafiek van $a^x$ en de transformaties van grafieken die je al kent, schets je snel de grafiek van een verwante functie.

\voorbeeld{}{Schets de grafiek van $g(x) = 2^{x + 1} - 3$. Geef het domein, het bereik en de asymptoot.\label{vb:h1-grafiek}}{%
\naast{%
\uitleg{De grafiek van $2^x$ schuift 1 naar links en 3 naar onder. De asymptoot schuift mee naar onder.}
\begin{stappenvrij}{4.2cm}{3.4cm}
(0, 1) \mapsto (-1, -2) & Het vaste punt. \\
\text{asymptoot } y = -3 & De $x$-as, 3 naar onder. \\
\dom g = \R & \\
\ber g = \left]-3, +\infty\right[ & Boven de asymptoot.
\end{stappenvrij}}{%
\begin{tikzpicture}[x=.5cm,y=.5cm,font=\scriptsize]
  \rooster{-4}{1}{2}{-3}{1}{3}
  \draw[grijs,dashed] (-4,-3) -- (2.4,-3);
  \kromme{-4}{1.32}{2^(\x+1)-3}
  \vol{(-1,-2)}
\end{tikzpicture}}}

\begin{oefeningen}
\oef{h1-k-grafiek}
\end{oefeningen}

\subsection{Exponentiële vergelijkingen}

Een exponentiële functie is overal stijgend of overal dalend. Twee verschillende exponenten kunnen dus nooit dezelfde macht geven.

\begin{theorie}
\begin{stelling}{Gelijke machten}
Voor $a > 0$ met $a \neq 1$ is \qquad $\formule{a^p = a^q \iff p = q}$
\end{stelling}
\begin{bewijs}
Is $p = q$, dan is zeker $a^p = a^q$. Omgekeerd: stel $p \neq q$, zeg $p < q$. Volgens de vorige stelling is dan $a^p < a^q$ (als $a > 1$) of $a^p > a^q$ (als $0 < a < 1$). In beide gevallen is $a^p \neq a^q$.
\end{bewijs}
\end{theorie}

De functie $a^x$ is dus injectief: elke macht hoort bij precies één exponent. Daarmee los je vergelijkingen op waarin beide leden een macht van hetzelfde grondtal zijn.

\voorbeeld{}{Los op: $2^{x + 1} = \dfrac18$.\label{vb:h1-vergelijking}}{%
\uitleg{Schrijf $\frac18$ als macht van 2. Twee machten van 2 zijn gelijk als hun exponenten gelijk zijn.}
\begin{stappenbreed}
2^{x + 1} = 2^{-3} & $\frac18 = \frac{1}{2^3} = 2^{-3}$. \\
x + 1 = -3 & Gelijke machten: exponenten gelijkstellen. \\
x = -4 & Controle: $2^{-3} = \frac18$.
\end{stappenbreed}}

Soms is geen van beide leden een macht van het grondtal dat je ziet staan. Zoek dan een grondtal waarvan beide grondtallen een macht zijn.

\voorbeeld{}{Los op: $4^{x + 1} = 8^{x}$.\label{vb:h1-vergelijking2}}{%
\uitleg{$4$ en $8$ zijn allebei machten van 2. Schrijf beide leden als macht van 2 met de regel voor een macht van een macht.}
\begin{stappenbreed}
\bigl(2^2\bigr)^{x + 1} = \bigl(2^3\bigr)^{x} & $4 = 2^2$ en $8 = 2^3$. \\
2^{2x + 2} = 2^{3x} & Exponenten vermenigvuldigen. \\
2x + 2 = 3x \iff x = 2 & Gelijke machten. Controle: $4^3 = 64 = 8^2$.
\end{stappenbreed}}

Het lukt niet altijd. Bij $2^x = 5$ is $5$ geen macht van $2$ met een exponent die je kent. Je weet wel dat $2^2 = 4 < 5 < 8 = 2^3$, dus omdat $2^x$ stijgt, ligt de oplossing tussen 2 en 3. Om ze precies te vinden heb je een logaritme nodig (hoofdstuk 2). En $2^x = -4$ heeft geen oplossing: een macht van 2 is nooit negatief.

\begin{werkwijze}
Schrijf beide leden als een macht van hetzelfde grondtal (gebruik de rekenregels). Stel dan de exponenten gelijk: $a^p = a^q \iff p = q$. Lukt het niet om hetzelfde grondtal te vinden, dan heb je een logaritme nodig (hoofdstuk 2).
\end{werkwijze}

\begin{oefeningen}
\oef{h1-vergelijking}

\oef{h1-k-tweekanten}
\end{oefeningen}

%========================================================================================
\section{Exponentiële groei en verval}

Een stad telt $1000$ inwoners, en dat aantal neemt elk jaar met $10$\,\% toe. Na een jaar komen er $100$ bij, het jaar daarna al $110$, want de nieuwe inwoners tellen mee:

\figuren{\renewcommand{\arraystretch}{1.2}%
\begin{tabular}{c|ccccc}
jaar $t$ & $0$ & $1$ & $2$ & $3$ & $4$ \\ \hline
inwoners & $1000$ & $1100$ & $1210$ & $1331$ & $1464{,}1$
\end{tabular}}

De toename is elk jaar anders, maar elk jaar vermenigvuldig je met dezelfde factor $1{,}1$. Na $t$ jaar is het aantal dus $1000 \cdot 1{,}1^t$.

\begin{theorie}
Een grootheid $N$ verandert \textbf{\textcolor{accent}{exponentieel}} als ze per tijdseenheid met dezelfde \textbf{\textcolor{accent}{groeifactor}} $a > 0$ vermenigvuldigd wordt:
\[\formule{N(t) = N_0 \cdot a^t}\]
met $N_0 = N(0)$ de beginwaarde. Is $a > 1$, dan spreken we van \textbf{\textcolor{accent}{groei}}, is $0 < a < 1$, van \textbf{\textcolor{accent}{verval}}.
\end{theorie}

\subsection{Groeifactor en procent}

Meestal krijg je geen groeifactor, maar een percentage. Hoe vind je de factor?

\begin{theorie}
\begin{stelling}{Groeifactor uit een percentage}
Neemt $N$ per tijdseenheid met $p$\,\% toe, dan is $\formule{a = 1 + \tfrac{p}{100}}$. \\
Neemt $N$ per tijdseenheid met $p$\,\% af, dan is $\formule{a = 1 - \tfrac{p}{100}}$.
\end{stelling}
\begin{bewijs}
Een toename met $p$\,\% van $N$ is $\frac{p}{100} \cdot N$. Daarna is er
\bewijsregel{N + \frac{p}{100} \cdot N = \Bigl(1 + \frac{p}{100}\Bigr) \cdot N .}
Bij een afname is er $N - \frac{p}{100} \cdot N = \bigl(1 - \frac{p}{100}\bigr) \cdot N$ over.
\end{bewijs}
\end{theorie}

Een vast percentage per tijdseenheid geeft dus exponentiële verandering. Omgekeerd lees je het percentage af uit de factor: $a = 1{,}07$ is $7$\,\% groei, $a = 0{,}93$ is $7$\,\% verval.

\voorbeeld{}{Je zet $2000$ euro op een rekening met $3$\,\% intrest per jaar. De intrest wordt elk jaar bijgeteld. Hoeveel staat er na 10 jaar op?\label{vb:h1-intrest}}{%
\uitleg{Elk jaar wordt het bedrag vermenigvuldigd met $1{,}03$: de intrest komt bij het kapitaal en levert het jaar erna zelf intrest op.}
\begin{stappenbreed}
K(t) = 2000 \cdot 1{,}03^t & Beginwaarde 2000, groeifactor $1 + \frac{3}{100}$. \\
K(10) = 2000 \cdot 1{,}03^{10} \approx 2687{,}83 & Na 10 jaar, in euro.
\end{stappenbreed}
Met enkelvoudige intrest (elk jaar $60$ euro) was dat maar $2600$ euro geweest.}

\naast{%
Het verschil met enkelvoudige intrest lijkt klein, maar het groeit. Bij enkelvoudige intrest komt er elk jaar hetzelfde \textit{bedrag} bij: dat is lineaire groei, met als grafiek een rechte. Bij samengestelde intrest wordt het bedrag elk jaar met dezelfde \textit{factor} vermenigvuldigd. Na 40 jaar staat er $2000 \cdot 1{,}03^{40} \approx 6524$ euro tegenover $2000 + 40 \cdot 60 = 4400$ euro.

Op lange termijn wint exponentiële groei het altijd van lineaire groei.}{%
\begin{tikzpicture}[x=.16cm,y=.55cm,font=\scriptsize]
  \draw[black!14,very thin,xstep=.8cm,ystep=.55cm] (0,0) grid (40,7);
  \draw[->,grijs] (0,0) -- (42.5,0) node[right,text=black] {$t$};
  \draw[->,grijs] (0,0) -- (0,7.5) node[above,text=black] {$K$ (in 1000 euro)};
  \foreach \t in {10,20,30,40}{\node[below,inner sep=1.5pt] at (\t,0) {$\t$};}
  \foreach \k in {1,...,7}{\node[left,inner sep=1.5pt] at (0,\k) {$\k$};}
  \kromme{0}{40}{2*1.03^(\x)}
  \draw[accent,thick,dashed] (0,2) -- (40,4.4);
  \node[above left,text=accent] at (38,6.2) {samengesteld};
  \node[below right,text=accent] at (24,3.4) {enkelvoudig};
\end{tikzpicture}}

\begin{oefeningen}
\oef{h1-k-intrest}
\end{oefeningen}

Bij verval moet je goed lezen: wat verdwijnt er, en wat blijft er over? De groeifactor is altijd wat er overblijft.

\voorbeeld{}{Een zwembad verliest elke dag $4$\,\% van zijn chloor. Er zit $1{,}5$ kg chloor in het water. Hoeveel is er na een week over, en met hoeveel procent is de hoeveelheid dan gedaald?\label{vb:h1-daling}}{%
\uitleg{Er verdwijnt $4$\,\% per dag, dus er blijft $96$\,\% over: de groeifactor is $0{,}96$.}
\begin{stappenbreed}
C(t) = 1{,}5 \cdot 0{,}96^t & In kg, $t$ in dagen. \\
C(7) = 1{,}5 \cdot 0{,}96^7 \approx 1{,}127 & Na een week, in kg. \\
0{,}96^7 \approx 0{,}751\,4 & Er blijft $75{,}1$\,\% over: een daling van $24{,}9$\,\%.
\end{stappenbreed}
De daling is niet $7 \cdot 4\,\% = 28$\,\%: elke dag verdwijnt $4$\,\% van een kleinere hoeveelheid.}

\begin{oefeningen}
\oef{h1-k-daling}
\end{oefeningen}

\subsection{Groeifactor voor een andere tijdseenheid}

Een groeifactor geldt per tijdseenheid. Wat gebeurt er over een langere of kortere periode?

\begin{theorie}
\begin{stelling}{Groeifactor over een periode}
Is $N(t) = N_0 \cdot a^t$, dan wordt $N$ in elke periode van $h$ tijdseenheden met $a^h$ vermenigvuldigd:
\[\formule{N(t + h) = N(t) \cdot a^h}.\]
\end{stelling}
\begin{bewijs}
Met de productregel voor machten is
\bewijsregel{N(t + h) = N_0 \cdot a^{t + h} = N_0 \cdot a^t \cdot a^h = N(t) \cdot a^h .}
\end{bewijs}
\end{theorie}

De factor $a^h$ hangt niet af van het moment $t$ waarop je begint: bij gelijke tijdsstappen hoort altijd dezelfde factor. Per 10 jaar is de factor $a^{10}$, per maand (een twaalfde jaar) is hij $a^{\frac{1}{12}}$.

\voorbeeld{}{Het aantal fietsers in een stad stijgt elk jaar met $6$\,\%. Met hoeveel procent stijgt het per maand?\label{vb:h1-tijdseenheid}}{%
\uitleg{Een maand is $\frac{1}{12}$ jaar. De groeifactor per maand is dus de jaarfactor tot de macht $\frac{1}{12}$, en niet $\frac{6}{12} = 0{,}5$\,\%.}
\begin{stappenbreed}
a = 1{,}06 & Groeifactor per jaar. \\
a^{\frac{1}{12}} = 1{,}06^{\frac{1}{12}} \approx 1{,}004\,87 & Groeifactor per maand. \\
1{,}004\,87 - 1 = 0{,}004\,87 & Ongeveer $0{,}49$\,\% per maand.
\end{stappenbreed}
Controle: $1{,}004\,87^{12} \approx 1{,}06$. Twaalf keer $0{,}5$\,\% was iets te veel geweest, want ook de groei groeit.}

\begin{oefeningen}
\oef{h1-groei}
\end{oefeningen}

\subsection{Verdubbelings- en halveringstijd}

Bij radioactief verval of bij medicijnen in het bloed geeft men meestal geen percentage, maar de tijd waarin de hoeveelheid halveert.

\begin{theorie}
De \textbf{\textcolor{accent}{verdubbelingstijd}} is de tijd waarin $N$ verdubbelt, de \textbf{\textcolor{accent}{halveringstijd}} de tijd waarin $N$ halveert.

\begin{stelling}{Model met verdubbelings- of halveringstijd}
Is $T$ de verdubbelingstijd, dan is $\formule{N(t) = N_0 \cdot 2^{\frac{t}{T}}}$. \\
Is $T$ de halveringstijd, dan is $\formule{N(t) = N_0 \cdot \Bigl(\frac12\Bigr)^{\frac{t}{T}}}$.
\end{stelling}
\begin{bewijs}
Neem de halveringstijd; voor de verdubbelingstijd vervang je $\frac12$ door $2$. Is $N(t) = N_0 \cdot a^t$, dan is $N(t + T) = N(t) \cdot a^T$, en dat moet $\frac12 N(t)$ zijn. Dus is $a^T = \frac12$, en met de regel voor een macht van een macht
\bewijsregel{a = \bigl(a^T\bigr)^{\frac1T} = \Bigl(\frac12\Bigr)^{\frac1T},}
zodat
\bewijsregel{N(t) = N_0 \cdot \Bigl(\Bigl(\frac12\Bigr)^{\frac1T}\Bigr)^t = N_0 \cdot \Bigl(\frac12\Bigr)^{\frac{t}{T}} .}
\end{bewijs}
\end{theorie}

De exponent $\frac{t}{T}$ telt hoeveel halveringstijden er voorbij zijn. Na één halveringstijd is de helft over, na twee een kwart, na drie een achtste. En omdat $a^T$ niet van $t$ afhangt, halveert de hoeveelheid in elke periode $T$, niet alleen in de eerste.

\voorbeeld{}{Cafeïne heeft in het lichaam een halveringstijd van ongeveer 5 uur. Iemand neemt 200 mg. Hoeveel is er na 12 uur over? Met hoeveel procent daalt de hoeveelheid per uur?\label{vb:h1-verval}}{%
\uitleg{Na $t$ uur zijn er $\frac{t}{5}$ halveringstijden voorbij. De groeifactor per uur is de factor $\frac12$ voor 5 uur, verdeeld over 5 gelijke stappen: $\bigl(\frac12\bigr)^{\frac15}$.}
\begin{stappenbreed}
C(t) = 200 \cdot \Bigl(\frac12\Bigr)^{\frac{t}{5}} & In mg, $t$ in uur. \\
C(12) = 200 \cdot 0{,}5^{2{,}4} \approx 37{,}9 & Na 12 uur, in mg. \\
a = 0{,}5^{\frac15} \approx 0{,}8706 & Groeifactor per uur. \\
1 - 0{,}8706 = 0{,}1294 & Ongeveer $12{,}9$\,\% minder per uur.
\end{stappenbreed}}

\begin{oefeningen}
\oef{h1-k-verval}
\end{oefeningen}

\needspace{12\baselineskip}
Met een verdubbelingstijd werk je op dezelfde manier.

\voorbeeld{}{Het aantal gebruikers van een nieuwe app verdubbelt elke 3 maanden. Nu zijn er $5000$. Hoeveel zijn het er na 10 maanden?\label{vb:h1-verdubbeling}}{%
\uitleg{Tien maanden zijn $\frac{10}{3}$ verdubbelingstijden. Dat is geen geheel getal, maar de formule werkt ook dan.}
\begin{stappenbreed}
G(t) = 5000 \cdot 2^{\frac{t}{3}} & $t$ in maanden. \\
G(10) = 5000 \cdot 2^{\frac{10}{3}} \approx 50\,400 & Ongeveer $50\,400$ gebruikers.
\end{stappenbreed}
Controle: na 9 maanden zijn er $5000 \cdot 2^3 = 40\,000$, na 12 maanden $80\,000$. De uitkomst ligt ertussen.}

\begin{werkwijze}
Zoek de beginwaarde $N_0$ en de groeifactor $a$ per tijdseenheid: uit een percentage, of uit een verdubbelings- of halveringstijd $T$ (dan is $N(t) = N_0 \cdot 2^{t/T}$ of $N_0 \cdot \bigl(\frac12\bigr)^{t/T}$). Een groeifactor voor een periode van $h$ tijdseenheden is $a^h$.
\end{werkwijze}

\begin{oefeningen}
\oef{h1-k-verdubbeltijd}
\end{oefeningen}

%========================================================================================
\section{Het getal e}

Wat gebeurt er als de intrest vaker wordt bijgeteld? Bij $100$\,\% per jaar, $n$ keer per jaar bijgeteld, krijg je telkens $\frac{100}{n}$\,\% intrest. Dan groeit $1$ euro in een jaar tot $\bigl(1 + \frac1n\bigr)^n$:

\figuren{\renewcommand{\arraystretch}{1.2}%
\begin{tabular}{c|ccccc}
$n$ & $1$ & $10$ & $100$ & $1000$ & $1\,000\,000$ \\ \hline
$\bigl(1 + \frac1n\bigr)^n$ & $2$ & $2{,}5937$ & $2{,}7048$ & $2{,}7169$ & $2{,}7182\,8$
\end{tabular}}

Het resultaat groeit niet onbeperkt, maar nadert een vast getal.

\begin{theorie}
Het \textbf{\textcolor{accent}{getal van Euler}} is het getal dat $\bigl(1 + \frac1n\bigr)^n$ nadert als $n$ onbeperkt groeit:
\[\formule{\e \approx 2{,}718\,28}.\]
De functie $\formule{f(x) = \e^x}$ heet de \textbf{\textcolor{accent}{natuurlijke exponentiële functie}}.
\end{theorie}

\naast{%
Het getal $\e$ is irrationaal, net als $\pi$. Omdat $2 < \e < 3$, ligt de grafiek van $\e^x$ rechts van de $y$-as tussen die van $2^x$ en $3^x$, en links ervan net omgekeerd.

Waarom juist dit grondtal \textit{natuurlijk} is, zie je bij de afgeleiden: $\e^x$ is de enige exponentiële functie die haar eigen afgeleide is.}{%
\begin{tikzpicture}[x=.8cm,y=.5cm,font=\scriptsize]
  \rooster{-2}{1}{2}{-1}{1}{6}
  \kromme[densely dotted]{-2}{2}{2^(\x)}
  \kromme{-2}{1.79}{exp(\x)}
  \kromme[dashed]{-2}{1.63}{3^(\x)}
  \vol{(0,1)}
  \node[right,text=accent] at (2,4) {$2^x$};
  \node[right,text=accent] at (1.8,5.9) {$\e^x$};
  \node[left,text=accent] at (1.6,6.2) {$3^x$};
\end{tikzpicture}}

\subsection{Continue groei}

Groeit een bedrag met een intrest $r$ per jaar (bijvoorbeeld $r = 0{,}05$ voor $5$\,\%), maar wordt de intrest $n$ keer per jaar bijgeteld, telkens met $\frac{r}{n}$, dan is de groeifactor per jaar $\bigl(1 + \frac{r}{n}\bigr)^n$. Wordt $n$ heel groot, dan spreken we van \textit{continue} groei. Schrijf $n = m \cdot r$, dan is
\[\Bigl(1 + \frac{r}{n}\Bigr)^n = \Bigl(1 + \frac{1}{m}\Bigr)^{m r} = \biggl(\Bigl(1 + \frac{1}{m}\Bigr)^{m}\biggr)^{r} .\]
Als $n$ onbeperkt groeit, doet $m$ dat ook, en nadert de binnenste macht het getal $\e$. De groeifactor per jaar nadert dus $\e^r$, en na $t$ jaar is het bedrag met $\e^{rt}$ vermenigvuldigd. Een volledig bewijs vraagt limieten; deze redenering toont waar de formule vandaan komt.

\begin{theorie}
Een grootheid $N$ groeit \textbf{\textcolor{accent}{continu}} met groeisnelheid $r$ als
\[\formule{N(t) = N_0 \cdot \e^{rt}}.\]
Bij $r > 0$ is dat groei, bij $r < 0$ verval.
\end{theorie}

\voorbeeld{}{Een bedrag van $1000$ euro groeit continu met $5$\,\% per jaar: $K(t) = 1000\,\e^{0{,}05\,t}$. Hoeveel is dat na 10 jaar? Vergelijk met $5$\,\% per jaar, jaarlijks bijgeteld.\label{vb:h1-e}}{%
\uitleg{Bereken beide met het rekentoestel (de toets $\e^x$).}
\begin{stappenbreed}
K(10) = 1000\,\e^{0{,}5} \approx 1648{,}72 & Continu. \\
1000 \cdot 1{,}05^{10} \approx 1628{,}89 & Jaarlijks: iets minder.
\end{stappenbreed}}

\begin{oefeningen}
\oef{h1-e}
\end{oefeningen}

Continue groei is ook gewoon exponentiële groei, maar met een groeifactor die je eerst moet uitrekenen.

\begin{theorie}
\begin{stelling}{Groeifactor bij continue groei}
Is $N(t) = N_0 \cdot \e^{rt}$, dan is de groeifactor per tijdseenheid \qquad $\formule{a = \e^r}$
\end{stelling}
\begin{bewijs}
Met de regel voor een macht van een macht is
\bewijsregel{N_0 \cdot \e^{rt} = N_0 \cdot \bigl(\e^r\bigr)^t .}
Dat is het model $N_0 \cdot a^t$ met $a = \e^r$.
\end{bewijs}
\end{theorie}

De groeisnelheid $r$ is dus niet het groeipercentage. Bij $r = 0{,}05$ is $a = \e^{0{,}05} \approx 1{,}0513$: per jaar komt er $5{,}13$\,\% bij, iets meer dan $5$\,\%. Dat zag je ook in het vorige voorbeeld.

\voorbeeld{}{Een radioactieve stof vervalt continu: $N(t) = 80\,\e^{-0{,}03\,t}$, met $N$ in gram en $t$ in jaar. Met hoeveel procent daalt de hoeveelheid per jaar?\label{vb:h1-continu}}{%
\uitleg{Schrijf het model met een groeifactor per jaar, en lees daaruit het percentage.}
\begin{stappenbreed}
a = \e^{-0{,}03} \approx 0{,}970\,4 & Groeifactor per jaar. \\
1 - 0{,}970\,4 = 0{,}029\,6 & Ongeveer $2{,}96$\,\% minder per jaar.
\end{stappenbreed}
Bij continu verval is het percentage per jaar iets kleiner dan $3$\,\%.}

\begin{oefeningen}
\oef{h1-k-e}
\end{oefeningen}

%========================================================================================
\opdrachtenkop

\begin{oefeningen}
\needspace{14\baselineskip}
\oef{h1-o1}

\needspace{12\baselineskip}
\oef{h1-o10}

\needspace{13\baselineskip}
\oef{h1-o2}

\needspace{12\baselineskip}
\oef{h1-o12}

\needspace{12\baselineskip}
\oef{h1-o3}

\oef{h1-o15}

\needspace{10\baselineskip}
\oef{h1-o13}

\needspace{16\baselineskip}
\oef{h1-o11}

\oef{h1-o14}

\needspace{12\baselineskip}
\oef{h1-o4}

\oef{h1-o18}

\oef{h1-o22}

\oef{h1-o16}

\oef{h1-o6}

\oef{h1-o21}

\oef{h1-o5}

\oef{h1-o19}

\oef{h1-o23}

\needspace{14\baselineskip}
\oef{h1-o17}

\oef{h1-o7}

\oef{h1-o20}

\needspace{10\baselineskip}
\oef{h1-o26}

\oef{h1-o8}

\oef{h1-o25}

\oef{h1-o27}

\oef{h1-o24}

\oef{h1-o9}
\end{oefeningen}

%----------------------------------------------------------------------------------------
% KOPPEL: verbanden met andere leerstof
%----------------------------------------------------------------------------------------
\needspace{20\baselineskip}
\opdrachtsectie{Koppel}

\lkr{Rijen, Kansrekenen en Integralen: kijk na wat al gezien is. Het pakket legt de volgorde niet vast.}

\begin{oefeningen}
\oef{h1-koppel-1}

\oef{h1-koppel-2}

\oef{h1-koppel-3}
\end{oefeningen}

%----------------------------------------------------------------------------------------
\opdrachtsectie{Test jezelf}

\begin{oefeningen}
\oef{h1-t1}

\oef{h1-t2}

\oef{h1-t3}

\oef{h1-t4}

\oef{h1-t5}

\oef{h1-t6}

\oef{h1-t7}
\end{oefeningen}

\herhaalschema{%
  \fv{oef:h1-t1}{vb:h1-macht}{oef:h1-machten}%
  \fv{oef:h1-t2}{vb:h1-vergelijking2}{oef:h1-k-tweekanten}%
  \fv{oef:h1-t3}{vb:h1-grafiek}{oef:h1-k-grafiek}%
  \fv{oef:h1-t4}{vb:h1-intrest}{oef:h1-k-intrest}%
  \fv{oef:h1-t5}{vb:h1-verval}{oef:h1-k-verval}%
  \fv{oef:h1-t6}{vb:h1-e}{oef:h1-k-e}%
  \fv{oef:h1-t7}{vb:h1-tijdseenheid}{oef:h1-groei}}

\ifzelfstudie
\section*{Maak je eigen samenvatting}
Sluit dit boek en schrijf in je schrift: wat $a^{m/n}$ en $a^{-r}$ betekenen, de rekenregels voor machten, de eigenschappen van de grafiek van $a^x$ voor $a > 1$ en voor $0 < a < 1$, het groeimodel $N_0 \cdot a^t$ met groeifactor, groeipercentage en halveringstijd, en de definitie van $\e$. Vergelijk daarna met de kaders.
\fi
